\documentclass[aspectratio=169,11pt]{beamer}

\usepackage[utf8]{inputenc}
\usepackage[T1]{fontenc}
\usepackage[spanish]{babel}
\usepackage{amsmath,amssymb}
\usepackage{booktabs}
\usepackage{graphicx}
\usepackage{tikz}

\definecolor{main}{RGB}{20,60,110}

\setbeamercolor{frametitle}{fg=white,bg=main}
\setbeamercolor{structure}{fg=main}
\setbeamercolor{block title}{bg=main,fg=white}
\setbeamercolor{block body}{bg=main!8}
\setbeamercolor{block title alerted}{bg=red!70!black,fg=white}
\setbeamercolor{block body alerted}{bg=red!8}
\setbeamertemplate{navigation symbols}{}
\setbeamertemplate{footline}{\hfill\scriptsize Newton Sketch \quad \insertframenumber/\inserttotalframenumber\hspace{1em}\vspace{0.5em}}

% \eqrow{Título}{ecuación}{explicación}
\newcommand{\eqrow}[3]{%
  \textcolor{main}{\textbf{#1}}\par\vspace{2pt}
  \begin{columns}[c]
    \begin{column}{0.55\textwidth}
      \centering $\displaystyle #2$
    \end{column}
    \begin{column}{0.42\textwidth}
      #3
    \end{column}
  \end{columns}
  \vspace{4pt}
}

% Espacio reservado para una imagen: \imgspace{alto}{descripción}
\newcommand{\imgspace}[2]{%
  \fbox{\parbox[c][#1][c]{0.9\linewidth}{\centering\small #2}}}

% Notación
\newcommand{\R}{\mathbb{R}}
\newcommand{\E}{\mathbb{E}}
\newcommand{\norm}[1]{\lVert #1\rVert}
\newcommand{\nablaf}{\nabla f}
\newcommand{\Hess}{\nabla^2 f}
\DeclareMathOperator{\nnz}{nnz}
\DeclareMathOperator{\diag}{diag}

\title{Newton Sketch}
\subtitle{Optimización de segundo orden con proyecciones aleatorias}
\author{Diego Arévalo y Yeferson Getial}
\institute{Optimización Numérica y Álgebra Lineal Numérica}
\date{}

\begin{document}

% ------------------------------------------------------------
\begin{frame}
  \titlepage
\end{frame}

% ============================================================
% 1. INTRODUCCIÓN Y MOTIVACIÓN
% ============================================================
\begin{frame}{Introducción}
  Queremos resolver un problema convexo, con $n \gg d$,
  \[
    \min_{x\in\R^d} f(x), \qquad f(x)=\sum_{i=1}^{n}\phi_i(a_i^\top x),
    \qquad A=\begin{bmatrix}a_1^\top\\ \vdots\\ a_n^\top\end{bmatrix}\in\R^{n\times d}.
  \]
  \vspace{-4pt}
  \begin{itemize}
    \item Regresión logística, mínimos cuadrados, programación con barreras.
    \pause
    \item[] \begin{alertblock}{Restricción de memoria}
      La matriz $A$, muy alta ($n\gg d$), puede no caber en memoria.
    \end{alertblock}
    \pause
    \item Nos gustaría un método con \alert{pocas iteraciones}, \alert{iteraciones baratas} y que \alert{use poca memoria}.
    \item Lo más importante: \alert{garantizar convergencia}.
  \end{itemize}
\end{frame}

\begin{frame}{Motivación: ¿por qué Newton y por qué es caro?}
  \begin{columns}[T]
    \begin{column}{0.48\textwidth}
      \textcolor{main}{\textbf{Primer orden (gradiente)}}
      \begin{itemize}
        \item Costo por iteración: $O(nd)$.
        \item Muchas iteraciones; sensible al condicionamiento.
      \end{itemize}
    \end{column}
    \begin{column}{0.48\textwidth}
      \textcolor{main}{\textbf{Segundo orden (Newton)}}
      \begin{itemize}
        \item Pocas iteraciones, \alert{invariante afín}.
        \item Costo por iteración: $O(nd^2 + d^3)$.
      \end{itemize}
    \end{column}
  \end{columns}
  \vspace{8pt}
  \pause
  \begin{block}{Pregunta}
    ¿Podemos conservar el comportamiento de Newton reduciendo la dimensión $n$ con aleatoriedad?
  \end{block}
\end{frame}

\begin{frame}{Índice}
  \begin{enumerate}
    \item \textbf{Motivación}
    \item \textbf{Funciones autoconcordantes}
    \item \textbf{Matrices de sketch} y \textbf{subspace embeddings}
    \item \textbf{Newton Sketch}
    \item \textbf{Experimentos, resultados y conclusiones}
  \end{enumerate}
\end{frame}

% ============================================================
% 2. FUNCIONES AUTOCONCORDANTES
% ============================================================
\begin{frame}{Funciones autoconcordantes}
  \begin{columns}[T]
    \begin{column}{0.58\textwidth}
      \begin{block}{Definición (una variable)}
        Una función convexa $f:I\subset\R\to\R$ de clase $C^3$ es \emph{autoconcordante} si, para todo $x\in I$,
        \[
          |f'''(x)|\;\le\;2\,f''(x)^{3/2}.
        \]
      \end{block}
      \pause
      \begin{itemize}
        \item Equivale a $\bigl|\tfrac{d}{dx}\,f''(x)^{-1/2}\bigr|\le 1$: la curvatura no cambia bruscamente respecto a sí misma.
        \item Ejemplo: $f(x)=-\log x$ ($f''=x^{-2}$, $|f'''|=2x^{-3}$, se cumple con igualdad).
        \item Más ejemplos: funciones lineales y cuadráticas.
      \end{itemize}
    \end{column}
    \begin{column}{0.38\textwidth}
      \centering
      % \includegraphics[width=\linewidth]{fig_autoconcordante_1d.pdf}
      \imgspace{5.5cm}{Imagen: gráfica de $-\log x$ (y su curvatura $f''$)}
    \end{column}
  \end{columns}
\end{frame}

\begin{frame}{Intuición: norma local y elipsoide de Dikin}
  \begin{columns}[T]
    \begin{column}{0.58\textwidth}
      \textbf{En varias variables:} $f$ es autoconcordante si $\varphi(t)=f(x+th)$ lo es para todo $x$ y $h$. Equivalentemente,
      \[
        \bigl|D^3 f(x)[h,h,h]\bigr| \le 2\,\bigl(h^\top \Hess(x)\,h\bigr)^{3/2}.
      \]
      \pause
      \begin{itemize}
        \item \textbf{Norma local:} $\norm{h}_x=\sqrt{h^\top\Hess(x)h}$.
        \item \textbf{Elipsoide de Dikin:} $\mathcal E_x=\{y:\norm{y-x}_x\le 1\}$ está contenido en el dominio de $f$.
        \item Dentro de $\mathcal E_x$ la Hessiana casi no cambia $\Rightarrow$ el modelo cuadrático de Newton es confiable.
      \end{itemize}
    \end{column}
    \begin{column}{0.38\textwidth}
      \centering
      % \includegraphics[width=\linewidth]{fig_dikin.pdf}
      \imgspace{3.2cm}{Imagen: elipsoides de Dikin en distintos puntos}
      \vspace{6pt}

      \begin{alertblock}{Dos fases}
        \small
        \textbf{Amortiguada:} $\lambda$ grande, paso $\mu<1$, descenso fijo.\\
        \textbf{Cuadrática:} $\lambda$ pequeño, el paso completo queda dentro de $\mathcal E_x$.
      \end{alertblock}
    \end{column}
  \end{columns}
\end{frame}



\begin{frame}{Newton en funciones autoconcordantes}
 
\end{frame}

% ============================================================
% 3. MATRICES DE SKETCH Y SUBSPACE EMBEDDING
% ============================================================
\begin{frame}{Matrices de sketch}
  \begin{columns}[c]
    \begin{column}{0.55\textwidth}
      Una matriz de \emph{sketch} $S\in\R^{m\times n}$, con $m\ll n$, es una proyección aleatoria que comprime filas:
      \[
        A\in\R^{n\times d}\ \longmapsto\ SA\in\R^{m\times d}.
      \]
      
    \end{column}
    \begin{column}{0.4\textwidth}
      \centering
      \begin{tikzpicture}[scale=0.8]
        \draw[fill=main!15] (0,0) rectangle (1,3);
        \node at (0.5,1.5) {$A$};
        \node at (-0.4,1.5) {$n$};
        \node at (0.5,3.3) {$d$};
        \draw[fill=main!35] (-3,1.0) rectangle (-1.2,1.0+0.9);
        \node at (-2.1,1.45) {$S$};
        \node at (-2.1,2.2) {$m\times n$};
        \node at (0,-0.5) {$\xrightarrow{\ \ SA\ \ }$};
      \end{tikzpicture}
    \end{column}
  \end{columns}

  \pause
  \begin{itemize}
        \item Se construye \alert{sin mirar los datos} (\emph{oblivious}) o adaptada a ellos.
        \item Debe preservar la información relevante de $A$.
        \item Se evalúa por costo de aplicar $SA$ y por tamaño $m$.
      \end{itemize}
\end{frame}

\begin{frame}{Subspace embedding}
  \begin{definition}
    $S\in\R^{m\times n}$ es un \emph{$\varepsilon$-subspace embedding} de $\operatorname{col}(A)$ si, para todo $y\in\R^d$,
    \[
      (1-\varepsilon)\norm{Ay}_2^2 \;\le\; \norm{SAy}_2^2 \;\le\; (1+\varepsilon)\norm{Ay}_2^2 .
    \]
  \end{definition}

  \pause
  
  \begin{itemize}
    \item Equivalente: $\norm{A^\top S^\top S A - A^\top A}_2\le\varepsilon\norm{A}_2^2$ (tras ortonormalizar).
    \item Consecuencia clave: $(SA)^\top(SA)\approx A^\top A$, es decir, una \alert{Hessiana aproximada} con relación espectral $1\pm\varepsilon$.
  \end{itemize}
\end{frame}

% ------------------------------------------------------------
\begin{frame}{Ejemplo 1: sketch Gaussiano}
  \eqrow{Definición}
    {S_{ij}\sim\mathcal{N}\!\left(0,\tfrac{1}{m}\right)\ \text{i.i.d.}}
    {Matriz densa $m\times n$.}
  \eqrow{Tamaño}
    {m=O\!\left(\dfrac{d+\log(1/\delta)}{\varepsilon^2}\right)}
    {Mínimo número de filas; garantía óptima.}
  

\end{frame}

\begin{frame}{Ejemplo 2: CountSketch}
  \begin{columns}[c]
    \begin{column}{0.5\textwidth}
      Cada columna de $S$ tiene \alert{un solo} elemento no nulo:
      \[
        S_{h(j),\,j}=\sigma_j,\quad h(j)\in[m],\ \sigma_j\in\{\pm1\}.
      \]
      \begin{itemize}
        \item $SA$: cada fila de $A$ se suma con signo en un \emph{bucket}.
        \item Costo: $O(\nnz(A))$.
        \item Requiere $m=O(d^2/(\varepsilon^2\delta))$.
      \end{itemize}
    \end{column}
    \begin{column}{0.45\textwidth}
      \centering
      \[
        S=\begin{bmatrix}
          0 & \!-1 & 0 & 1 & 0\\
          1 & 0 & 0 & 0 & -1\\
          0 & 0 & 1 & 0 & 0
        \end{bmatrix}
      \]
      \footnotesize Ejemplo con $m=3$, $n=5$.
    \end{column}
  \end{columns}
\end{frame}


\begin{frame}{Comparación de sketches}
  \centering
  \begin{tabular}{lccc}
    \toprule
    & \textbf{Gaussiana} & \textbf{CountSketch} & \textbf{SRHT} \\
    \midrule
    Filas $m$ & $O(d/\varepsilon^2)$ & $O(d^2/\varepsilon^2)$ & $\tilde O(d/\varepsilon^2)$ \\
    Costo de $SA$ & $O(mnd)$ & $O(\nnz(A))$ & $O(nd\log m)$ \\
    Estructura & densa & dispersa & rápida (FWHT) \\
    \bottomrule
  \end{tabular}
  \vspace{10pt}

  \begin{block}{Compromiso}
    Menos filas $\Leftrightarrow$ aplicar $S$ más caro. Newton Sketch debe elegir según $n$, $d$ y $\nnz(A)$.
  \end{block}
\end{frame}

% ============================================================
% 4. NEWTON SKETCH
% ============================================================
\begin{frame}{Estructura de la Hessiana}
  Para $f(x)=\sum_i\phi_i(a_i^\top x)$:
  \\
  \eqrow{Gradiente}
    {\nablaf(x)=A^\top v(x),\quad v_i=\phi_i'(a_i^\top x)}
    {Costo $O(nd)$.}
    \pause
  \eqrow{Hessiana}
    {\Hess(x)=A^\top D(x)\,A,\quad D=\diag\bigl(\phi_i''(a_i^\top x)\bigr)}
    {$D\succeq 0$ diagonal.}
    \pause
  \eqrow{Raíz cuadrada}
    {\Hess(x)=\bigl(D^{1/2}A\bigr)^{\!\top}\bigl(D^{1/2}A\bigr)}
    {La raíz $D^{1/2}A$ se obtiene \alert{gratis} escalando filas.}
\end{frame}

\begin{frame}{Newton Sketch: el algoritmo}
  \textbf{Idea:} sustituir la raíz exacta por una versión comprimida.
  \vspace{4pt}
  \begin{alertblock}{Newton Sketch}
    \[
      \widehat H_t=\bigl(S_t D_t^{1/2}A\bigr)^{\!\top}\bigl(S_t D_t^{1/2}A\bigr),
      \qquad
      x_{t+1}=x_t-\mu_t\,\widehat H_t^{-1}\nablaf(x_t).
    \]
  \end{alertblock}
  \pause
  \begin{enumerate}
    \item Calcular $\nablaf(x_t)$ y $D_t^{1/2}A$.
    \item Generar $S_t\in\R^{m\times n}$ \alert{independiente} en cada iteración.
    \item Formar $S_tD_t^{1/2}A$ y resolver el sistema $\widehat H_t\,\Delta=-\nablaf(x_t)$.
    \item Búsqueda lineal sobre $f$ y actualizar $x_{t+1}=x_t+\mu_t\Delta$.
  \end{enumerate}
  \textbf{Costo:} $O(\text{costo}(SD^{1/2}A)+md^2+d^3)$ en lugar de $O(nd^2+d^3)$.
\end{frame}

\begin{frame}{El problema: la diagonal y la estructura}
  \begin{itemize}
    \item $D_t$ cambia en cada iteración $\Rightarrow$ hay que recalcular $D_t^{1/2}A$ y \alert{re-sketchear}.
    \item Escalar filas cuesta $O(\nnz(A))$ y se hace ``gratis'', pero \dots
    \item Con $S$ Gaussiana densa el producto domina: $O(mnd)$.
    \item Con CountSketch se conserva $O(\nnz(A))$, pero $m=O(d^2)$ crece.
    \item Si $D$ es casi singular (muchos $\phi_i''\approx 0$), el sketch pierde rango efectivo.
  \end{itemize}
  \begin{block}{Idea}
    Elegir $S$ según la estructura: Gaussiana/SRHT cuando $m$ pequeño importa; CountSketch cuando $A$ es dispersa.
  \end{block}
  % TODO: ajustar esta diapositiva a la interpretación exacta de "diagonal" que usen en su trabajo.
\end{frame}

% ============================================================
% 5. EJEMPLO PUNTUAL
% ============================================================
\begin{frame}{Ejemplo puntual: una iteración}
  \textbf{Regresión logística} con $n=6$, $d=2$ (reemplazar con sus datos reales).
  \begin{columns}[T]
    \begin{column}{0.5\textwidth}
      \begin{enumerate}
        \item $x_0=(0,0)^\top$, entonces $D_0=\tfrac14 I$.
        \item $\nablaf(x_0)=A^\top v_0=\;?$
        \item $D_0^{1/2}A$ y sketch $S\in\R^{3\times 6}$.
        \item $\widehat H_0=\;?$ \quad vs.\quad $H_0=\;?$
      \end{enumerate}
    \end{column}
    \begin{column}{0.46\textwidth}
      \centering
      \begin{tabular}{lcc}
        \toprule
        & Newton & Newton Sketch \\
        \midrule
        $\Delta$ & $?$ & $?$ \\
        $f(x_1)$ & $?$ & $?$ \\
        $\lambda(x_1)$ & $?$ & $?$ \\
        \bottomrule
      \end{tabular}
    \end{column}
  \end{columns}
\end{frame}

% ============================================================
% 6. DEMOSTRACIÓN
% ============================================================
\begin{frame}{Resultado principal}
  
  \centering\small Convergencia \alert{lineal-cuadrática}: tasa lineal controlada por $\varepsilon$, no por el número de condición.
\end{frame}

\begin{frame}{Bosquejo de la demostración}
  
\end{frame}

% ============================================================
% 7. EXPERIMENTOS Y RESULTADOS
% ============================================================
\begin{frame}{Experimentos: configuración}
  \begin{itemize}
    \item \textbf{Problemas:} regresión logística, mínimos cuadrados, \dots
    \item \textbf{Datos:} sintéticos ($n=10^5$--$10^6$, $d=100$--$500$) y reales (\texttt{nombre del dataset}).
    \item \textbf{Métodos comparados:} Newton exacto, gradiente descendente, L-BFGS, Newton Sketch (Gaussiana / CountSketch / SRHT).
    \item \textbf{Parámetros:} $m=c\,d$ con $c\in\{2,4,8\}$; criterio de parada $\lambda(x)^2/2\le 10^{-8}$.
    \item \textbf{Métricas:} error $f(x_t)-f^\star$ vs.\ iteraciones y vs.\ tiempo.
  \end{itemize}
\end{frame}

\begin{frame}{Resultados}
  \begin{columns}[c]
    \begin{column}{0.52\textwidth}
      \centering
      % \includegraphics[width=\textwidth]{fig_convergencia.pdf}
      \imgspace{4.2cm}{Figura: $f(x_t)-f^\star$ vs.\ iteraciones}
    \end{column}
    \begin{column}{0.44\textwidth}
      \centering\small
      \begin{tabular}{lcc}
        \toprule
        Método & Iter. & Tiempo (s) \\
        \midrule
        Newton & $?$ & $?$ \\
        NS Gaussiana & $?$ & $?$ \\
        NS CountSketch & $?$ & $?$ \\
        \bottomrule
      \end{tabular}
    \end{column}
  \end{columns}
\end{frame}

\begin{frame}{Resultados: efecto del tamaño del sketch}
  \centering
  % \includegraphics[width=0.7\textwidth]{fig_m.pdf}
  \fbox{\parbox[c][5cm][c]{0.7\textwidth}{\centering Figura: tasa de convergencia en función de $m/d$}}
  \vspace{6pt}

  \small Verifica que la tasa lineal mejora al aumentar $m$, ya que $\varepsilon\sim\sqrt{d/m}$.
\end{frame}

% ============================================================
% 8. CONCLUSIONES Y REFERENCIAS
% ============================================================
\begin{frame}{Conclusiones}
  \begin{itemize}
    \item \dots % completar
    \item \dots % completar
    \item \dots % completar
  \end{itemize}
\end{frame}

\begin{frame}[allowframebreaks]{Referencias}
  \footnotesize
  \begin{thebibliography}{9}
    \bibitem{pw2017} M.~Pilanci and M.~J.~Wainwright, \emph{Newton Sketch: A Near Linear-Time Optimization Algorithm with Linear-Quadratic Convergence}, SIAM J. Optim., 27(1), 2017.
    \bibitem{nn1994} Y.~Nesterov and A.~Nemirovskii, \emph{Interior-Point Polynomial Algorithms in Convex Programming}, SIAM, 1994.
    \bibitem{bv2004} S.~Boyd and L.~Vandenberghe, \emph{Convex Optimization}, Cambridge University Press, 2004.
    \bibitem{w2014} D.~P.~Woodruff, \emph{Sketching as a Tool for Numerical Linear Algebra}, Found.\ Trends Theor.\ Comput.\ Sci., 10(1--2), 2014.
    \bibitem{cw2013} K.~L.~Clarkson and D.~P.~Woodruff, \emph{Low Rank Approximation and Regression in Input Sparsity Time}, STOC, 2013.
  \end{thebibliography}
\end{frame}

\end{document}